\section{Products of binary monomials}\label{sec:binary}

The method of Sections~\ref{sec:coordinates}--\ref{sec:tests} extends
from $x_i^{d_i-1}y_i$ to binary monomials $x_i^{a_i}y_i^{b_i}$ with
$a_i\ge b_i\ge1$. Let $\mathbf a=(a_1,\ldots,a_k)$ and
$\mathbf b=(b_1,\ldots,b_k)$ be integer vectors with $a_i\ge b_i\ge1$,
put $d_i=a_i+b_i$, and consider
\begin{equation}\label{eq:bm-target}
F_{\mathbf a,\mathbf b}(x,y)=\prod_{i=1}^k x_i^{a_i}y_i^{b_i}.
\end{equation}
Its partially symmetric rank $R_{\bd}(F_{\mathbf a,\mathbf b})$ is the
least $r$ for which $F_{\mathbf a,\mathbf b}$ admits an expression of
the form \eqref{eq:decomposition}, with the blocks $(x_i,y_i)$ and the
degrees $d_i$ fixed. Interchanging $x_i$ and $y_i$ preserves the form;
for arbitrary positive exponents the rank is obtained by replacing
$(a_i,b_i)$ with $(\max\{a_i,b_i\},\min\{a_i,b_i\})$. When every
$b_i=1$, $F_{\mathbf a,\mathbf b}$ is the polynomial \eqref{eq:target}.

\begin{theorem}[Products of binary monomials]\label{thm:binary}
Let $k\ge1$, let $a_i\ge b_i\ge1$ be integers, and put $d_i=a_i+b_i$.
Over $\C$,
\begin{equation}\label{eq:bm-rank}
R_{\bd}\!\left(\prod_{i=1}^k x_i^{a_i}y_i^{b_i}\right)
=N(\mathbf a,\mathbf b)
:=\prod_{i=1}^k(b_i+1)+\sum_{i=1}^k(a_i-b_i)\prod_{j\ne i}(b_j+1),
\end{equation}
where empty products equal $1$.
\end{theorem}

The lower bound holds for every finite decomposition
\eqref{eq:decomposition} of $F_{\mathbf a,\mathbf b}$, without
restrictions on its complex coefficients or its support. The upper
bound is realized by the explicit identity \eqref{eq:bm-character}.
The theorem concerns the exact complex rank with the symmetric blocks
fixed. It makes no assertion about real rank, border rank, cactus rank,
unrestricted nonsymmetric rank, or rank after grouping modes from
different blocks, and the lower-bound argument uses only finite
decompositions. Pure-power factors, with $b_i=0$, are excluded; see
\Cref{rem:bm-pure}.

The lower bound is first proved for abstract product moment conditions
in \Cref{thm:pm}, and \Cref{prop:bm-reduction} reduces an arbitrary
decomposition to such conditions. The upper bound is the
finite-character decomposition of \Cref{prop:bm-upper}. In this
section, $m_i(e)$ denotes general moment values as in \Cref{thm:pm};
the values \eqref{eq:mi} are a special case.

\subsection{A product-moment lower bound}
Throughout this subsection, $b_i\ge1$ and $d_i\ge2b_i$ are integers
for $i\in[k]$, and
\begin{equation}\label{eq:bm-parameters}
\delta_i:=d_i-2b_i,\qquad
q_i:=\prod_{j\ne i}(b_j+1),\qquad
N:=\prod_{i=1}^k(b_i+1)+\sum_{i=1}^k\delta_iq_i .
\end{equation}
When every $b_j=1$, these are $\delta_i=d_i-2$, $q_i=q=2^{k-1}$ and
$N=M$; see \eqref{eq:Mcount}.

\begin{theorem}[Product-moment lower bound]\label{thm:pm}
Let $k\ge1$, let $b_i\ge1$ and $d_i\ge2b_i$ be integers, and let $N$
be as in \eqref{eq:bm-parameters}. For each $i$, let
$m_i(0),\ldots,m_i(d_i)\in\C$ satisfy
\begin{equation}\label{eq:bm-truncation}
m_i(e)=0\quad(b_i<e\le d_i),
\qquad
\beta_i:=m_i(b_i)\ne0 .
\end{equation}
Let $z_1,\ldots,z_r\in\C^k$ and nonzero $w_1,\ldots,w_r\in\C$ satisfy
\begin{equation}\label{eq:bm-moments}
\sum_{\ell=1}^r w_\ell\prod_{i=1}^k z_{i\ell}^{e_i}
=\prod_{i=1}^k m_i(e_i),
\qquad 0\le e_i\le d_i .
\end{equation}
Then $r\ge N$. The points may repeat, and any of their coordinates may
be zero.
\end{theorem}

Fix data as in \Cref{thm:pm} for the rest of this subsection. Let
$L_i$ be the linear functional on univariate polynomials of degree at
most $d_i$ determined by $L_i(Z_i^e)=m_i(e)$, and let $L$ be the
weighted evaluation \eqref{eq:L}. The proof of \Cref{lem:generalized}
uses only linearity, so the lemma holds verbatim with
\eqref{eq:bm-moments} in place of \eqref{eq:moments}. In particular,
the product rule \eqref{eq:product} holds whenever $\deg P_i\le d_i$
for every $i$; all Gram computations below use \eqref{eq:bm-moments}
only in this form.

Let $t\in\C^k$ be admissible, that is, $1+t_iz_{i\ell}\ne0$ for every
$i,\ell$, and let $h_i$ and $H_t$ be as in \eqref{eq:Ht}. For rational
functions $\phi,\psi$ defined at every point $z_\ell$, put
\begin{equation}\label{eq:bm-pairing}
\langle\phi,\psi\rangle_t
:=\sum_{\ell=1}^r w_\ell H_t(z_\ell)\phi(z_\ell)\psi(z_\ell).
\end{equation}
This is a symmetric bilinear form; no positivity is involved. The Gram
matrix of a finite list of such functions is the matrix \eqref{eq:Gab}
of that list, so \Cref{lem:gram,lem:dimension} apply to it in place of
$\phi_1,\ldots,\phi_M$. The identities below hold for every admissible
$t$; equivalently, they hold over the field $\C(t_1,\ldots,t_k)$ with
indeterminate parameters, in which every $1+t_iz_{i\ell}$ is a nonzero
polynomial, also when $z_{i\ell}=0$.

Write $\mathcal I:=\prod_{j=1}^k\{0,\ldots,b_j\}$ and
$\mathcal I_i:=\prod_{j\ne i}\{0,\ldots,b_j\}$. Empty products,
including empty Kronecker products of matrices, equal $1$ (the
$1\times1$ matrix $(1)$); when $k=1$, $\mathcal I_1$ consists of the
empty tuple. The basic list is
\begin{equation}\label{eq:bm-U}
\mathscr U:=(f_\iota)_{\iota\in\mathcal I},
\qquad
f_\iota(Z):=\prod_{j=1}^k\frac{Z_j^{\iota_j}}{h_j(Z_j)^{b_j}},
\end{equation}
ordered lexicographically. For each $i$ with $\delta_i>0$, the
additional list $\mathscr V_i$ consists of the functions
\begin{equation}\label{eq:bm-V}
g_{i,s,\iota}(Z):=\frac{Z_i^s}{h_i(Z_i)^s}
\prod_{j\ne i}\frac{Z_j^{\iota_j}}{h_j(Z_j)^{b_j}},
\qquad
b_i+1\le s\le d_i-b_i,\quad \iota\in\mathcal I_i,
\end{equation}
ordered first by increasing $s$ and then lexicographically by $\iota$;
if $\delta_i=0$, then $\mathscr V_i$ is empty. Thus $\mathscr U$ has
$\prod_j(b_j+1)$ members and $\mathscr V_i$ has $\delta_iq_i$ members,
$N$ in total. When every $b_j=1$, $\mathscr U$ and $\mathscr V_i$ are
the families $\cB_0$ and $\cB_i$ of \eqref{eq:B0}--\eqref{eq:Bi}.

For $0\le e,e'\le b_i$, define the $(b_i+1)\times(b_i+1)$ matrix
$G_i(t_i)$ by
\begin{equation}\label{eq:bm-Gi}
[G_i(t_i)]_{e,e'}:=L_i\!\left(Z_i^{e+e'}h_i(Z_i)^{\delta_i}\right).
\end{equation}

\begin{lemma}[The basic block]\label{lem:bm-basic}
For every admissible $t$, the Gram matrix of $\mathscr U$ is
$\bigotimes_{i=1}^kG_i(t_i)$, and
\begin{equation}\label{eq:bm-Gdet}
\det G_i(t_i)=(-1)^{b_i(b_i+1)/2}\beta_i^{b_i+1}\ne0 .
\end{equation}
In particular, the Gram matrix of $\mathscr U$ is invertible for every
admissible $t$.
\end{lemma}
\begin{proof}
The polynomial in \eqref{eq:bm-Gi} has degree at most
$2b_i+\delta_i=d_i$. If $e+e'>b_i$, each of its monomials has degree
between $b_i+1$ and $d_i$, so the entry vanishes by
\eqref{eq:bm-truncation}. If $e+e'=b_i$, only the lowest term
survives, and the entry equals $\beta_i$. Reversing the order of the
columns therefore gives an upper triangular matrix whose diagonal
entries all equal $\beta_i$; the reversal of $b_i+1$ columns has sign
$(-1)^{b_i(b_i+1)/2}$, which proves \eqref{eq:bm-Gdet}. Finally, for
$\iota,\jmath\in\mathcal I$, multiplication by $H_t$ cancels the
denominators, and the product rule gives
\begin{equation*}
\langle f_\iota,f_\jmath\rangle_t
=L\!\left(\prod_{i=1}^k Z_i^{\iota_i+\jmath_i}h_i(Z_i)^{\delta_i}\right)
=\prod_{i=1}^k[G_i(t_i)]_{\iota_i,\jmath_i},
\end{equation*}
since the $i$th factor has degree at most $d_i$. This is the
$(\iota,\jmath)$ entry of the Kronecker product.
\end{proof}

\begin{lemma}[Orthogonality of the lists]\label{lem:bm-orth}
For every admissible $t$, the lists
$\mathscr U,\mathscr V_1,\ldots,\mathscr V_k$ are mutually orthogonal
for \eqref{eq:bm-pairing}.
\end{lemma}
\begin{proof}
Pair $g_{i,s,\iota}$ with $f_\jmath$. After multiplication by $H_t$,
the factor in coordinate $i$ is
\[
Z_i^{s+\jmath_i}h_i(Z_i)^{d_i-s-b_i}.
\]
The exponent of $h_i$ is nonnegative because $s\le d_i-b_i$. Every
monomial has degree at least $s>b_i$ and at most
$d_i-b_i+\jmath_i\le d_i$, so its $L_i$-value is zero. Every other
coordinate factor is $Z_j^{\iota_j+\jmath_j}h_j(Z_j)^{\delta_j}$, of
degree at most $d_j$. The product rule gives
$\langle g_{i,s,\iota},f_\jmath\rangle_t=0$.

Now let $i\ne j$, and pair $g_{i,s,\iota}\in\mathscr V_i$ with
$g_{j,s',\jmath}\in\mathscr V_j$. The coordinate-$i$ factor has the
same form $Z_i^{s+\jmath_i}h_i(Z_i)^{d_i-s-b_i}$, with $\jmath_i$
supplied by the second function, and again has zero $L_i$-value. The
coordinate-$j$ factor $Z_j^{s'+\iota_j}h_j(Z_j)^{d_j-s'-b_j}$ has a
nonnegative exponent on $h_j$ and degree at most $d_j$. All remaining
coordinate factors have the form used in \Cref{lem:bm-basic}. Thus
every use of the product rule stays inside the degree box, and
$\langle g_{i,s,\iota},g_{j,s',\jmath}\rangle_t=0$.
\end{proof}

When every $b_j=1$, these are the computations of
Appendix~\ref{app:entries}. It remains to treat the Gram matrix of a
nonempty $\mathscr V_i$. Define the $q_i\times q_i$ matrix $H_i(t)$,
with rows and columns indexed by $\iota,\jmath\in\mathcal I_i$, by
\begin{equation}\label{eq:bm-Hi}
[H_i(t)]_{\iota,\jmath}
:=\sum_{\ell=1}^r w_\ell\frac{z_{i\ell}^{d_i+1}}{1+t_iz_{i\ell}}
\prod_{j\ne i}z_{j\ell}^{\iota_j+\jmath_j}(1+t_jz_{j\ell})^{\delta_j}.
\end{equation}
When every $b_j=1$, this is the matrix \eqref{eq:Hi}.

\begin{lemma}[The anti-triangular blocks]\label{lem:bm-anti}
Let $\delta_i>0$, and group $\mathscr V_i$ by $s$. For every admissible
$t$, the $(s,s')$ block of the Gram matrix of $\mathscr V_i$ is zero if
$s+s'\le d_i$ and equals $H_i(t)$ if $s+s'=d_i+1$. Consequently, up to
the sign of a block-column permutation, the Gram determinant of
$\mathscr V_i$ equals $\det H_i(t)^{\delta_i}$.
\end{lemma}
\begin{proof}
The $(s,s')$ block has entries
\begin{equation}\label{eq:bm-block}
\sum_{\ell=1}^r w_\ell z_{i\ell}^{s+s'}(1+t_iz_{i\ell})^{d_i-s-s'}
\prod_{j\ne i}z_{j\ell}^{\iota_j+\jmath_j}(1+t_jz_{j\ell})^{\delta_j}.
\end{equation}
If $s+s'\le d_i$, this is the $L$-value of a polynomial in $\cP_{\bd}$
whose coordinate-$i$ degrees all exceed $b_i$; hence the whole block
is zero. If $s+s'=d_i+1$, the block is exactly $H_i(t)$.

Write $s=b_i+1+\alpha$ with $0\le\alpha<\delta_i$. After reversing the
order of the block columns, the block column with index $\alpha'$
corresponds to $s'=d_i-b_i-\alpha'$, so
\[
s+s'=d_i+1+\alpha-\alpha'.
\]
Every block above the diagonal ($\alpha'>\alpha$) is therefore zero,
and every diagonal block is $H_i(t)$. The resulting block lower
triangular matrix has determinant $\det H_i(t)^{\delta_i}$, as
asserted. The blocks with $s+s'\ge d_i+2$ lie below the diagonal and
are not needed.
\end{proof}

The next identity shows that $H_i$ is nonsingular as a rational
matrix. Its form includes zero coordinates of the points.

\begin{lemma}[A finite division identity]\label{lem:bm-division}
For integers $b\ge1$ and $d\ge2b$, the identity
\begin{equation}\label{eq:bm-division}
\frac{z^{d+1}}{1+tz}
=\sum_{e=b}^{d}(-1)^{d-e}\frac{z^{e}}{t^{d+1-e}}
+(-1)^{d+1-b}\frac{z^{b}}{t^{d+1-b}(1+tz)}
\end{equation}
holds whenever $t\ne0$ and $1+tz\ne0$. For each fixed $z\ne0$, the
final term is $O(t^{-(d-b+2)})$ as $t\to\infty$; for $z=0$, every term
of \eqref{eq:bm-division} vanishes identically. For $b=1$, the
substitution $n=d-e$ turns \eqref{eq:bm-division} into
\eqref{eq:division}.
\end{lemma}
\begin{proof}
Multiply \eqref{eq:bm-division} by $(1+tz)t^{d+1-b}$. Consecutive terms
of the finite sum cancel, leaving $t^{d+1-b}z^{d+1}$.
If $z\ne0$, the additional denominator $1+tz$ gives one further power
of $t^{-1}$ at infinity. If $z=0$, the factor $z^b$ vanishes because
$b\ge1$, as do all the other numerators. Thus neither the identity nor
the remainder statement requires division by $z$.
\end{proof}

\begin{lemma}[Nonsingularity of the additional blocks]\label{lem:bm-Hi}
Let $\delta_i>0$. Then $\det H_i(t)$ is a nonzero element of
$\C(t_1,\ldots,t_k)$. More precisely, as a Laurent expansion in
$t_i^{-1}$ with coefficients in $\C(t_j:j\ne i)$,
\begin{equation}\label{eq:bm-leading}
H_i(t)=(-1)^{d_i-b_i}\beta_i\,t_i^{-(d_i-b_i+1)}
\bigotimes_{j\ne i}G_j(t_j)+O\!\left(t_i^{-(d_i-b_i+2)}\right),
\end{equation}
and the leading coefficient matrix is invertible.
\end{lemma}
\begin{proof}
Treat the parameters $t_j$, $j\ne i$, as elements of the coefficient
field $\C(t_j:j\ne i)$. For $\iota,\jmath\in\mathcal I_i$, put
$A_{\iota,\jmath}(Z_{-i}):=
\prod_{j\ne i}Z_j^{\iota_j+\jmath_j}h_j(Z_j)^{\delta_j}$;
its degree in each $Z_j$ is at most $d_j$. Applying
\eqref{eq:bm-division} with $d=d_i$ and $b=b_i$ separately at each
point in \eqref{eq:bm-Hi} gives
\begin{equation}\label{eq:bm-expanded}
[H_i(t)]_{\iota,\jmath}
=\sum_{e=b_i}^{d_i}\frac{(-1)^{d_i-e}}{t_i^{d_i+1-e}}\,
L\!\left(Z_i^{e}A_{\iota,\jmath}\right)
+\frac{(-1)^{d_i+1-b_i}}{t_i^{d_i+1-b_i}}
\sum_{\ell=1}^r w_\ell
\frac{z_{i\ell}^{b_i}A_{\iota,\jmath}(z_\ell)}{1+t_iz_{i\ell}}.
\end{equation}
The sum over the points is finite and its terms with $z_{i\ell}=0$
vanish identically, so by \Cref{lem:bm-division} the last term is
$O(t_i^{-(d_i-b_i+2)})$ entrywise. For $b_i<e\le d_i$, the product
rule and \eqref{eq:bm-truncation} give $L(Z_i^eA_{\iota,\jmath})=0$.
For $e=b_i$, the product rule gives
\[
L\!\left(Z_i^{b_i}A_{\iota,\jmath}\right)
=\beta_i\prod_{j\ne i}
L_j\!\left(Z_j^{\iota_j+\jmath_j}h_j(Z_j)^{\delta_j}\right)
=\beta_i\left[\bigotimes_{j\ne i}G_j(t_j)\right]_{\iota,\jmath}.
\]
This proves \eqref{eq:bm-leading}, which may equivalently be read as a
formal Laurent expansion in $t_i^{-1}$. By \Cref{lem:bm-basic}, the
leading coefficient matrix
$C_i:=(-1)^{d_i-b_i}\beta_i\bigotimes_{j\ne i}G_j(t_j)$ is invertible.
The coefficient of $t_i^{-(d_i-b_i+1)q_i}$ in $\det H_i$ is therefore
$\det C_i\ne0$, so $\det H_i$ is not the zero rational function. Each
coordinate $i$ is treated separately; no common limiting path for the
parameters is required. When every $b_j=1$, \eqref{eq:bm-leading} is
\eqref{eq:laurent}.
\end{proof}

\begin{proof}[Proof of \Cref{thm:pm}]
For each $i$, put $D_i(t_i):=\prod_{\ell=1}^r(1+t_iz_{i\ell})^{q_i}$.
If $\delta_i>0$, the entries of $H_i$ have denominators dividing
$\prod_\ell(1+t_iz_{i\ell})$, so $D_i\det H_i\in\C[t]$; it is a
nonzero polynomial by \Cref{lem:bm-Hi}. Hence the polynomial
\begin{equation}\label{eq:bm-avoidance}
\left(\prod_{i,\ell}(1+t_iz_{i\ell})\right)
\prod_{i:\,\delta_i>0}D_i(t_i)\det H_i(t)
\end{equation}
is nonzero. Since $\C$ is infinite, it is nonzero at some finite
$t\in\C^k$; fix such a $t$. Then $t$ is admissible, so all the
functions \eqref{eq:bm-U}--\eqref{eq:bm-V} and the pairing
\eqref{eq:bm-pairing} are defined, and $\det H_i(t)\ne0$ whenever
$\delta_i>0$. The Gram matrix of $\mathscr U$ is invertible by
\Cref{lem:bm-basic}, and each nonempty $\mathscr V_i$ has an
invertible Gram matrix by \Cref{lem:bm-anti}. By \Cref{lem:bm-orth},
the Gram matrix of the combined list
$\mathscr U,\mathscr V_1,\ldots,\mathscr V_k$ is block diagonal with
these blocks, hence invertible. \Cref{lem:gram,lem:dimension}, applied
to these $N$ functions, show that their evaluation vectors are
linearly independent in $\C^r$, so $r\ge N$.

No point was discarded or required to be distinct or to have nonzero
coordinates. If a moment formula \eqref{eq:bm-moments} initially
contains zero weights, those terms may first be deleted.
\end{proof}

As in Section~\ref{sec:tests}, the count rests on the invertibility of
every diagonal Gram block, established before dimensions are added;
no positivity of the weights is used, and the auxiliary parameters are
evaluated at one finite point.

\subsection{From decompositions to product moments}

\begin{proposition}[Reduction to moments]\label{prop:bm-reduction}
Let $a_i\ge b_i\ge1$ and $d_i=a_i+b_i$. Every decomposition
\eqref{eq:decomposition} of $F_{\mathbf a,\mathbf b}$ has at least
$N(\mathbf a,\mathbf b)$ summands.
\end{proposition}
\begin{proof}
To avoid a clash with the exponents $a_i,b_i$, write the decomposition
as
\begin{equation}\label{eq:bm-decomposition}
F_{\mathbf a,\mathbf b}(x,y)=\sum_{\ell=1}^r c_\ell
\prod_{i=1}^k(u_{i\ell}x_i+v_{i\ell}y_i)^{d_i},
\end{equation}
that is, \eqref{eq:decomposition} with $(a_{i\ell},b_{i\ell})$ renamed
$(u_{i\ell},v_{i\ell})$. Delete its zero summands; this does not
increase $r$, and afterwards every $c_\ell\ne0$ and every pair
$(u_{i\ell},v_{i\ell})\ne(0,0)$. As in the proof of
\Cref{lem:coordinates}, choose $\gamma_i\in\C$ with
$u_{i\ell}+\gamma_iv_{i\ell}\ne0$ for every $\ell$: a pair with
$v_{i\ell}\ne0$ excludes only $-u_{i\ell}/v_{i\ell}$, and a pair with
$v_{i\ell}=0$ excludes nothing. Set
\begin{equation}\label{eq:bm-zw}
z_{i\ell}=\frac{v_{i\ell}}{u_{i\ell}+\gamma_iv_{i\ell}},
\qquad
w_\ell=c_\ell\prod_{i=1}^k(u_{i\ell}+\gamma_iv_{i\ell})^{d_i},
\end{equation}
as in \eqref{eq:zw}. All weights are nonzero. The points may repeat,
and $z_{i\ell}$ may be zero.

For $0\le e_i\le d_i$, apply the differential operator
\[
\prod_{i=1}^k(\partial_{x_i}+\gamma_i\partial_{y_i})^{d_i-e_i}
\partial_{y_i}^{e_i}
\]
to \eqref{eq:bm-decomposition} and divide by $\prod_id_i!$. Since
$(\partial_{x_i}+\gamma_i\partial_{y_i})^{d_i-e_i}\partial_{y_i}^{e_i}
(u_{i\ell}x_i+v_{i\ell}y_i)^{d_i}
=d_i!\,(u_{i\ell}+\gamma_iv_{i\ell})^{d_i-e_i}v_{i\ell}^{e_i}$,
the result on the $\ell$th summand is $w_\ell\prod_iz_{i\ell}^{e_i}$.
On the target \eqref{eq:bm-target} the result factors over the blocks,
and the factor of the $i$th block is
\begin{equation}\label{eq:bm-mi}
m_i(e):=\frac{1}{d_i!}
(\partial_{x_i}+\gamma_i\partial_{y_i})^{d_i-e}\partial_{y_i}^{e}
\bigl(x_i^{a_i}y_i^{b_i}\bigr)
=\begin{cases}
\dfrac{(d_i-e)!\,b_i!}{(b_i-e)!\,d_i!}\,\gamma_i^{b_i-e},
&0\le e\le b_i,\\[8pt]
0,&b_i<e\le d_i.
\end{cases}
\end{equation}
Indeed, when $e\le b_i$, differentiation by $\partial_{y_i}^{e}$ first
contributes $b_i!/(b_i-e)!$. In the binomial expansion of
$(\partial_{x_i}+\gamma_i\partial_{y_i})^{d_i-e}$, only the term with
$b_i-e$ derivatives in $y_i$ and $a_i$ derivatives in $x_i$ survives;
it contributes $\binom{d_i-e}{b_i-e}\gamma_i^{b_i-e}a_i!\,(b_i-e)!$,
and the product of these factors divided by $d_i!$ is
\eqref{eq:bm-mi}. For $e>b_i$, the first differentiation already
vanishes.

Consequently the moment conditions \eqref{eq:bm-moments} hold
throughout the entire box $0\le e_i\le d_i$, with $m_i(e)$ given by
\eqref{eq:bm-mi}, and
\[
\beta_i=m_i(b_i)=\frac{a_i!\,b_i!}{d_i!}
=\binom{d_i}{b_i}^{-1}\ne0 .
\]
Because $a_i\ge b_i$, we have $d_i\ge2b_i$ and
$\delta_i=d_i-2b_i=a_i-b_i$, so the number $N$ of
\eqref{eq:bm-parameters} is $N(\mathbf a,\mathbf b)$.
\Cref{thm:pm} gives $r\ge N(\mathbf a,\mathbf b)$.
\end{proof}

The parameters $\gamma_i$ and the auxiliary parameters $t$ in
\Cref{thm:pm} serve different purposes. The first make an arbitrary
finite decomposition affine; the second make the Gram blocks
nondegenerate. Both choices are made after the decomposition is fixed,
and both follow from explicit finite nonvanishing conditions. Neither
is a genericity assumption on the decomposition.

\subsection{A finite-character decomposition}
In this subsection $a_i\ge b_i\ge1$, $d_i=a_i+b_i$,
$\delta_i=a_i-b_i$ and $N=N(\mathbf a,\mathbf b)$. Let
$\delta=(\delta_1,\ldots,\delta_k)^{\mathsf T}$,
$\mathbf 1=(1,\ldots,1)^{\mathsf T}$, and let
$\mathbf e_1,\ldots,\mathbf e_k$ be the standard basis vectors of
$\mathbb Z^k$. The construction uses a lattice whose nonzero vectors
avoid the box of possible exponent differences, followed by character
orthogonality on the corresponding finite quotient.

\begin{lemma}[A lattice avoiding the coefficient box]\label{lem:bm-lattice}
Let
\[
\Delta:=\diag(b_1+1,\ldots,b_k+1)+\delta\mathbf 1^{\mathsf T},
\]
and let $\Lambda:=\Delta\mathbb Z^k$ be the lattice generated by the
columns of $\Delta$. Then $\det\Delta=N$, so $[\mathbb Z^k:\Lambda]=N$,
and
\begin{equation}\label{eq:bm-box}
\Lambda\cap\prod_{i=1}^k[-a_i,b_i]=\{0\}.
\end{equation}
\end{lemma}
\begin{proof}
The $j$th column of $\Delta$ is $(b_j+1)\mathbf e_j+\delta$. Expand the
determinant multilinearly in these columns. Any term that chooses
$\delta$ from two columns vanishes. The term choosing no such column is
$\prod_i(b_i+1)$, and the term choosing it only from column $i$ is
$\delta_i\prod_{j\ne i}(b_j+1)$. Hence
\[
\det\Delta=\prod_{i=1}^k(b_i+1)
+\sum_{i=1}^k\delta_i\prod_{j\ne i}(b_j+1)=N>0.
\]
The index of the column lattice of a nonsingular integer matrix is the
absolute value of its determinant: integer row and column operations
preserve both quantities and reduce the matrix to diagonal form, where
the claim is immediate. Thus $[\mathbb Z^k:\Lambda]=N$.

Suppose that $\lambda=\Delta\mu$ lies in the box in \eqref{eq:bm-box},
with $\mu\in\mathbb Z^k$, and put $\sigma=\sum_i\mu_i$. Since $\Lambda$
is generated by the columns,
\begin{equation}\label{eq:bm-lattice-coordinates}
\lambda_i=(b_i+1)\mu_i+\delta_i\sigma .
\end{equation}
If $\sigma>0$, then $\lambda_i\le b_i$ and $\delta_i\ge0$ imply
$(b_i+1)\mu_i\le b_i$ for every $i$. Thus every $\mu_i\le0$, contrary
to $\sigma>0$. If $\sigma<0$, then $-\sigma-1\ge0$; since
$a_i=b_i+\delta_i$, the lower bound $\lambda_i\ge-a_i$ gives
\[
(b_i+1)\mu_i\ge-(b_i+1)+1+(-\sigma-1)\delta_i\ge-(b_i+1)+1 .
\]
Integrality now implies $\mu_i\ge0$ for every $i$, contrary to
$\sigma<0$. Finally, if $\sigma=0$, the upper bounds again imply
$\mu_i\le0$ for every $i$, and their zero sum forces $\mu=0$. This
proves \eqref{eq:bm-box}, including the cases in which some $\delta_i$
vanish.
\end{proof}

\begin{lemma}[Finite character orthogonality]\label{lem:bm-characters}
Let $\Gamma$ be a finite abelian group, and let
$\widehat\Gamma=\operatorname{Hom}(\Gamma,\C^\times)$. Then
$|\widehat\Gamma|=|\Gamma|$ and, for every $\xi\in\Gamma$,
\begin{equation}\label{eq:bm-character-sum}
\sum_{\chi\in\widehat\Gamma}\chi(\xi)=
\begin{cases}
|\Gamma|,&\xi=0,\\
0,&\xi\ne0.
\end{cases}
\end{equation}
\end{lemma}
\begin{proof}
Write $\Gamma$ as a product of cyclic groups
$\mathbb Z/\nu_1\mathbb Z\times\cdots\times\mathbb Z/\nu_\kappa\mathbb Z$.
A character is determined by independently choosing a $\nu_j$th root
of unity as the image of the $j$th cyclic generator. There are
therefore $\prod_j\nu_j=|\Gamma|$ characters. If $\xi$ has coordinates
$\xi_j$, the sum in \eqref{eq:bm-character-sum} factors as
\[
\prod_{j=1}^{\kappa}\left(\sum_{n=0}^{\nu_j-1}
\exp\!\left(\frac{2\pi\mathrm i\,n\xi_j}{\nu_j}\right)\right).
\]
Each geometric sum equals $\nu_j$ when $\xi_j=0$ in
$\mathbb Z/\nu_j\mathbb Z$ and equals zero otherwise, which proves the
formula.
\end{proof}

\begin{proposition}[An explicit decomposition]\label{prop:bm-upper}
Let $\Lambda$ be the lattice of \Cref{lem:bm-lattice}, put
$\Gamma:=\mathbb Z^k/\Lambda$, and for $\chi\in\widehat\Gamma$ put
$\theta_i(\chi):=\chi(\mathbf e_i+\Lambda)$. Then
\begin{equation}\label{eq:bm-character}
F_{\mathbf a,\mathbf b}(x,y)
=\frac{1}{N\prod_{i=1}^k\binom{d_i}{a_i}}
\sum_{\chi\in\widehat\Gamma}
\left(\prod_{i=1}^k\theta_i(\chi)^{-a_i}\right)
\prod_{i=1}^k\bigl(\theta_i(\chi)x_i+y_i\bigr)^{d_i}.
\end{equation}
This is a decomposition of the form \eqref{eq:decomposition} with $N$
nonzero summands. In particular,
$R_{\bd}(F_{\mathbf a,\mathbf b})\le N(\mathbf a,\mathbf b)$.
\end{proposition}
\begin{proof}
Each $\theta_i(\chi)$ is a root of unity, and in particular nonzero.
Let $\mathbf c=(c_1,\ldots,c_k)$ satisfy $0\le c_i\le d_i$. Since
$\prod_i\theta_i(\chi)^{c_i-a_i}=\chi(\mathbf c-\mathbf a+\Lambda)$, the
coefficient of $\prod_ix_i^{c_i}y_i^{d_i-c_i}$ on the right of
\eqref{eq:bm-character} is
\begin{equation}\label{eq:bm-coefficient}
\frac{\prod_i\binom{d_i}{c_i}}{N\prod_i\binom{d_i}{a_i}}
\sum_{\chi\in\widehat\Gamma}\chi(\mathbf c-\mathbf a+\Lambda).
\end{equation}
Since $d_i-a_i=b_i$, the vector $\mathbf c-\mathbf a$ lies in
$\prod_i[-a_i,b_i]$. By \Cref{lem:bm-lattice}, it belongs to
$\Lambda$ precisely when $\mathbf c=\mathbf a$. As $|\Gamma|=N$,
\Cref{lem:bm-characters} therefore makes \eqref{eq:bm-coefficient}
zero for $\mathbf c\ne\mathbf a$ and equal to one for
$\mathbf c=\mathbf a$. This
verifies every coefficient of \eqref{eq:bm-character}.

There are $|\widehat\Gamma|=N$ summands, and every coefficient in
\eqref{eq:bm-character} is nonzero. The tuples
$(\theta_1(\chi),\ldots,\theta_k(\chi))$ are also distinct: the classes
$\mathbf e_i+\Lambda$ generate $\Gamma$, so two characters with the
same $\theta_i$ for every $i$ coincide. Moreover, the product
$\prod_i(\theta_ix_i+y_i)^{d_i}$ has coefficient one at
$\prod_iy_i^{d_i}$ and coefficient $d_i\theta_i$ at
$x_iy_i^{d_i-1}\prod_{j\ne i}y_j^{d_j}$; thus distinct tuples give
nonproportional products. The identity is consequently an $N$-term
decomposition over $\C$ with the blocks fixed.
\end{proof}

\begin{proof}[Proof of \Cref{thm:binary}]
\Cref{prop:bm-reduction} bounds every decomposition below by
$N(\mathbf a,\mathbf b)$, and \Cref{prop:bm-upper} supplies a
decomposition with exactly this many summands.
\end{proof}

\subsection{Specializations}

\begin{remark}[The case $b_i=1$]\label{rem:bm-W}
Take $b_i=1$ and $a_i=d_i-1$, which is allowed for every $d_i\ge2$.
Then $\delta_i=d_i-2$ and
$N(\mathbf a,\mathbf b)=2^k+\sum_i(d_i-2)2^{k-1}=M$. The moments
\eqref{eq:bm-mi} become \eqref{eq:mi}, with $\beta_i=1/d_i$; the lists
$\mathscr U$ and $\mathscr V_i$ are $\cB_0$ and $\cB_i$; and
$G_i$, $H_i$ and \eqref{eq:bm-leading} reduce to \eqref{eq:Gi},
\eqref{eq:Hi} and \eqref{eq:laurent}. Thus \Cref{thm:binary} contains
\Cref{thm:main}, and gives
$R_{\bd}\bigl(\prod_ix_i^{d_i-1}y_i\bigr)=2^{k-1}(\sum_id_i-2k+2)$
for all $k\ge1$ and $d_i\ge2$. In particular,
\Cref{prop:bm-upper} gives an explicit $M$-term decomposition for all
$k\ge1$ and $d_i\ge2$, independently of
\cite[Theorem~1.1]{CCS2025}.
\end{remark}

For one block, \eqref{eq:bm-rank} gives $R_{(a+b)}(x^ay^b)=a+1$ for
$a\ge b\ge1$, the classical rank of a binary monomial
\cite[Section~4.3 of the arXiv version]{BBGOV2021}. For two blocks
it becomes
\begin{align*}
N(\mathbf a,\mathbf b)
&=(b_1+1)(b_2+1)+(a_1-b_1)(b_2+1)+(a_2-b_2)(b_1+1)\\
&=(a_1+1)(a_2+1)-(a_1-b_1)(a_2-b_2).
\end{align*}
For two binary-monomial blocks, Ga{\l}{\fontencoding{T1}\selectfont\k{a}}zka constructed a reduced
apolar set giving
\[
R_{\bd}\bigl(x_1^{a_1}y_1^{b_1}x_2^{a_2}y_2^{b_2}\bigr)
\le(a_1+1)(b_2+1)+(b_1+1)(a_2+1)-(b_1+1)(b_2+1),
\qquad a_i\ge b_i\ge1,
\]
and obtained lower bounds and exact special cases
\cite[Theorem~1.5 and Section~7.1 of the arXiv version]{Galazka2023}.
The binomial equations in that construction describe a finite set in
the two-dimensional torus. For $k=2$ the support of the decomposition
\eqref{eq:bm-character} is exactly the point set of
Ga{\l}{\fontencoding{T1}\selectfont\k{a}}zka's construction
\cite[Section~7.1 of the arXiv version]{Galazka2023}. The lower bound
\[
R_{\bd}\bigl(x_1^{a_1}y_1^{b_1}x_2^{a_2}y_2^{b_2}\bigr)
\ge\max\{(a_1+1)(b_2+1),(b_1+1)(a_2+1)\}
\]
also appears in
\cite[Proposition~4.3 of the arXiv version]{BBGOV2021} (stated there
with the roles of $x_i$ and $y_i$ exchanged). Behr's
unpublished manuscript \cite{Behr2026} gives the exact two-block
formula
\[
R_{\bd}\bigl(x_1^{a_1}y_1^{b_1}x_2^{a_2}y_2^{b_2}\bigr)
=(a_1+1)(a_2+1)-(a_1-b_1)(a_2-b_2)
\qquad(a_i\ge b_i\ge1)
\]
in Corollary~5.4, including the value $2(d_1+d_2-2)$ for
$x_1^{d_1-1}y_1x_2^{d_2-1}y_2$ with $d_1,d_2\ge2$ in Theorem~1.1.
Thus the two-block specialization is already present in that work; its
section on three or more factors leaves that case open.

When all $a_i=b_i$, the lists $\mathscr V_i$ are absent, the formula
reduces to $R_{\bd}\bigl(\prod_ix_i^{b_i}y_i^{b_i}\bigr)=\prod_i(b_i+1)$,
and the lattice matrix $\Delta$ is diagonal.

\begin{example}[Ranks $20$ and $54$]\label{ex:bm-54}
For the cubic triple $x_1^2y_1x_2^2y_2x_3^2y_3$ of \Cref{ex:rank20},
$b_i=1$ and $\delta_i=1$, so $N=2^3+3\cdot2^2=20$. The matrix
$\Delta$ has diagonal entries $3$ and off-diagonal entries $1$, and
\eqref{eq:bm-character} is a sum over the $20$ characters of
$\mathbb Z^3/\Delta\mathbb Z^3$. For $\mathbf a=(3,3,3)$ and
$\mathbf b=(2,2,2)$, we have $\delta_i=1$, so
\[
R_{(5,5,5)}\bigl(x_1^3y_1^2x_2^3y_2^2x_3^3y_3^2\bigr)=3^3+3\cdot3^2=54 .
\]
Here $\Delta$ has diagonal entries $4$ and off-diagonal entries $1$,
and \eqref{eq:bm-character} provides the $54$ summands.
\end{example}

\begin{remark}[Pure-power factors]\label{rem:bm-pure}
The hypothesis $b_i\ge1$ is essential to \Cref{thm:binary} as stated.
In the lower bound it is used where the remainder in
\eqref{eq:bm-division} is evaluated at a zero coordinate. The formula
would already be false for a single pure power $x^a$ with $b=0$ and
$a\ge1$: its rank is $1$, whereas the right-hand side of
\eqref{eq:bm-rank} would be $a+1$. Such factors are excluded from the
statement.
\end{remark}
